% !TeX root = ../main.tex
\section{引言}\label{sec:intro}
本项目比较两种室内温度校正方案。所有数据均为教学合成数据，
只用于练习论文排版，不能据此声称算法具有实际科研效果。
\emph{可复现的排版}要求正文、数据与引用能够一起交付。
数据版本为 \texttt{sensor\_v2}，抽查比例为 20\%。
\footnote{本实验不收集真实受试者数据。}
方法见第~\ref{sec:method} 节，结果与局限见第~\ref{sec:results} 节。
\subsection{研究问题}
比较 Baseline 和 Calibrated 在六个时刻的误差。
\begin{itemize}
  \item 保留原始观测；
  \item 图、表和正文使用同一组数据。
\end{itemize}
\begin{enumerate}
  \item 读取数据；\item 计算误差；\item 复核结果。
\end{enumerate}
\begin{description}
  \item[真实值] 本教学数据中给定的参照温度。
  \item[预测值] 两种方法给出的温度。
\end{description}
